\documentclass[11pt, a4paper, oneside]{article}
\usepackage[ngerman]{babel}
\usepackage{worksheet}

\begin{document}
	\author{L. Bung}
	\title{Webentwicklung: Frameworks}
	\subject{SAE}
	\class{E3FI}
	\maketitle


	\task[\timeLimit{$\sim$ 10 min.}]{Warum Frameworks?}

	In den letzten Stunden haben Sie Webseiten mit reinem HTML, CSS und JavaScript erstellt.
	In der Praxis wird jedoch selten alles ``von Hand'' geschrieben.

	\begin{enumerate}[label=\alph*)]
		\item Notieren Sie in Partnerarbeit mindestens drei Probleme, die bei großen Webprojekten mit reinem HTML, CSS und JavaScript entstehen können.
		\item Sammeln Sie Namen von Frameworks oder Bibliotheken aus dem Bereich Webentwicklung, die Sie kennen oder schon einmal gehört haben.
		\item Melden Sie sich, wenn Sie eines dieser Frameworks in Ihrem Ausbildungsbetrieb selbst verwenden.
		Daraus werden die Gruppen für Aufgabe 2 gebildet.
	\end{enumerate}

    \lines[6cm]


	\task[\timeLimit{$\sim$ 30 min.}]{Mini-Präsentationen vorbereiten}

	Die Klasse teilt sich in zwei Gruppen auf.
	Am Ende dieser Phase soll jede Gruppe (bzw. jedes Team) eine Mini-Präsentation von höchstens fünf Minuten vorbereitet haben.
	Alle Präsentationen orientieren sich an den folgenden Leitfragen:

	\begin{itemize}
		\item Wie heißt das Framework und wer entwickelt bzw. pflegt es?
		\item Um welche Art handelt es sich (z.\,B. CSS-Framework, Frontend-Framework, Bibliothek, Backend-Framework)?
		\item Welches Problem löst das Framework und wofür wird es typischerweise eingesetzt?
		\item Wie sieht ein minimales Beispiel aus?
        Präsentieren Sie einen kleinen Codeausschnitt.
		\item Welche Vor- und Nachteile hat das Framework?
	\end{itemize}

	\subsection*{Gruppe 2A: Framework-Profis}

	Sie verwenden in Ihrem Betrieb bereits ein Web-Framework (z.\,B. React, Angular, Vue.js, Svelte, jQuery, Next.js, \dots).
	Bilden Sie Teams von zwei bis drei Personen, die dasselbe Framework nutzen.

	\begin{enumerate}[label=\alph*)]
		\item Bereiten Sie eine Mini-Präsentation zu Ihrem Framework anhand der Leitfragen vor.
		\item Ergänzen Sie, wie und in welchem Projekt Sie das Framework im Betrieb einsetzen, ohne dabei vertrauliche Firmeninformationen preiszugeben.
		\item Zeigen Sie nach Möglichkeit ein kleines, selbst geschriebenes Codebeispiel live.
	\end{enumerate}

	\subsection*{Gruppe 2B: Bootstrap-Team}

	Sie haben im Betrieb wenig oder keine Berührung mit Webentwicklung oder zugehörigen Frameworks.
	Sie arbeiten sich in das weit verbreitete CSS-Framework \emph{Bootstrap} ein und stellen es anschließend der Klasse vor.
	Die offizielle Dokumentation finden Sie unter \url{https://getbootstrap.com}.

	\begin{enumerate}[label=\alph*)]
		\item Lesen Sie auf der Bootstrap-Webseite den Abschnitt ``Get started''.
		Erstellen Sie eine neue HTML-Datei und binden Sie Bootstrap über ein CDN ein, indem Sie im \texttt{head}-Bereich folgende Zeile ergänzen:
\begin{lstlisting}[language=html, breaklines=true, numbers=none, frame=none]
<link rel="stylesheet" href="https://cdn.jsdelivr.net/npm/bootstrap@5.3.3/dist/css/bootstrap.min.css">
\end{lstlisting}
		\item Bereiten Sie eine Mini-Präsentation zu Bootstrap vor, während Sie die folgenden Teilaufgaben bearbeiten.
		\item Informieren Sie sich über das Grid-System von Bootstrap.
		Erstellen Sie mit den Klassen \texttt{container}, \texttt{row} und \texttt{col-md-4} ein Layout mit drei Spalten, die auf einem Smartphone untereinander angezeigt werden.
		Überprüfen Sie das Verhalten mit den Entwicklertools Ihres Browsers.
		\item Kopieren Sie Ihr HTML-Formular aus den vergangenen Stunden in die Datei.
		Gestalten Sie es mithilfe der Bootstrap-Klassen \texttt{form-label}, \texttt{form-control} und \texttt{btn btn-primary} neu.
		\item Ergänzen Sie eine weitere Bootstrap-Komponente Ihrer Wahl, z.\,B. eine \texttt{navbar} oder eine \texttt{card}.
	\end{enumerate}


	\pagebreak

	\task[\timeLimit{$\sim$ 30 min.}]{Framework-Präsentationen}

	Die Teams stellen ihre Frameworks nacheinander vor.
	Halten Sie während der Präsentationen die wichtigsten Informationen in der folgenden Tabelle fest.

	\begin{center}
		%\renewcommand{\arraystretch}{3}
		\begin{tabular}{|p{2.8cm}|p{3.2cm}|p{3.2cm}|p{4.2cm}|}
			\hline
			\textbf{Framework} & \textbf{Art / Einsatzzweck} & \textbf{Besonderheit(en)} & \textbf{Vor- und Nachteile} \\
			\hline
			\hline
			Bootstrap\vspace{2.5cm} & & & \\
			\hline
			\phantom{x}\vspace{2.5cm} & & & \\
			\hline
			\phantom{x}\vspace{2.5cm} & & & \\
			\hline
			\phantom{x}\vspace{2.5cm} & & & \\
			\hline
			\phantom{x}\vspace{2.5cm} & & & \\
			\hline
		\end{tabular}
	\end{center}


	\pagebreak

	\task[\timeLimit{$\sim$ 15 min.}]{Das passende Werkzeug wählen}

	Nicht jedes Framework passt zu jedem Projekt.
	Entscheiden Sie in Partnerarbeit, welches der vorgestellten Frameworks (oder gar keines) Sie in den folgenden Situationen einsetzen würden, und begründen Sie Ihre Wahl.

	\begin{enumerate}[label=\alph*)]
		\item Ein kleiner Handwerksbetrieb möchte eine einfache Webseite mit Kontaktformular, die auf allen Geräten gut aussieht.
		\item Ein Start-up entwickelt eine Single-Page-Application mit vielen interaktiven Elementen, die sich ständig aktualisieren.
		\item Ein großes Unternehmen plant eine langlebige Unternehmensanwendung, an der mehrere Teams mit TypeScript und festen Strukturvorgaben arbeiten.
		\item Sie sollen ein internes Tool zur Erfassung von Urlaubsanträgen bauen, das schnell fertig sein muss und optisch nicht besonders anspruchsvoll ist.
	\end{enumerate}

	\lines[7cm]


	\task{Zusatzaufgabe}

	Bootstrap lässt sich mit einem JavaScript-Framework kombinieren.
	Recherchieren Sie, wie Bootstrap z.\,B. in React (Stichwort: \emph{React-Bootstrap}) oder Vue.js genutzt werden kann.
	Erklären Sie, warum die direkte Einbindung der Bootstrap-JavaScript-Datei in solchen Projekten zu Problemen führen kann.

\end{document}
